\begin{resumo}
Este trabalho analisa a sustentabilidade fiscal dos estados brasileiros no período de 2002 a 2023, com foco na reação do resultado primário ao endividamento e no papel dos contratos de refinanciamento da Lei n\textsuperscript{o}~9.496/1997. Adapta-se ao contexto subnacional brasileiro a estratégia empírica de Abubakar, McCausland e Theodossiou (2025), estimando-se por mínimos quadrados em dois estágios com efeitos fixos uma função em que o resultado primário/RCL responde à dívida/RCL defasada e à sua interação com o teto contratual. Os resultados corroboram a condição de sustentabilidade de Bohn, ao atestarem que, em média, os estados elevam o resultado primário em resposta ao incremento do endividamento. A intensidade dessa reação, todavia, é heterogênea conforme o desenho contratual: tetos mais permissivos atenuam o ajuste fiscal. O trabalho contribui para deslocar o foco da simples existência de regras fiscais para o desenho institucional dessas regras, explorando a heterogeneidade dos contratos de refinanciamento como elemento relevante para a sustentabilidade fiscal subnacional.
\vspace{\onelineskip}

\noindent
\textbf{Palavras-chave}: sustentabilidade fiscal; reação fiscal; regras fiscais; dívida estadual; endividamento subnacional\end{resumo}

\begin{resumo}[Abstract]
\begin{otherlanguage*}{english}
This study analyzes the fiscal sustainability of Brazilian states over the period 2002--2023, focusing on the reaction of the primary balance to indebtedness and on the role of the refinancing contracts under Law No.~9.496/1997. The empirical strategy of Abubakar, McCausland and Theodossiou (2025) is adapted to the Brazilian subnational context, estimating by two-stage least squares with fixed effects a function in which the primary balance relative to current net revenue (RCL) responds to lagged consolidated net debt relative to RCL and to its interaction with the contractual ceiling. The results corroborate Bohn's sustainability condition, attesting that, on average, states raise the primary balance in response to increased indebtedness. The intensity of this reaction, however, is heterogeneous according to the contractual design: more permissive ceilings attenuate the fiscal adjustment. The study contributes to shifting the focus from the mere existence of fiscal rules to the institutional design of these rules, exploring the heterogeneity of refinancing contracts as a relevant element for subnational fiscal sustainability.
\vspace{\onelineskip}

\noindent
\textbf{Keywords}: fiscal sustainability; state debt; fiscal rules; fiscal reaction function; subnational indebtedness\end{otherlanguage*}
\end{resumo}